\Chapter{ARTICLE 1~: ENERGY-EFFICIENT AND NEAR REAL-TIME ROUTING ALGORITHM FOR DISASTER MONITORING IN WIRELESS SENSOR NETWORKS}\label{sec:Theme1}

\vspace{0.5cm}
\noindent\textbf{Auteurs~:} Georges Parfait Djimefo Kapen, Ranwa Al Mallah et Samuel Pierre.

\noindent\textbf{Revue~:} IEEE Transactions on Green Communications and Networking.

\noindent\textbf{Statut~:} Soumis.

\vspace{0.3cm}
\noindent\textit{Ce chapitre pr\'{e}sente le premier article de cette th\`{e}se. Il propose GAPF (Graph-Attentive Policy Fusion), un cadre de m\'{e}ta-apprentissage l\'{e}ger qui fusionne des politiques d'experts MARL (apprentissage par renforcement multi-agent) gel\'{e}es \`{a} travers un encodeur \`{a} r\'{e}seau convolutif sur graphe (GCN) et un contr\^{o}leur de s\'{e}lection contextuelle d'algorithme (CAS). GAPF encode la topologie des capteurs en un plongement de graphe compact et s\'{e}lectionne l'expert le plus performant avec seulement 1\,473 param\`{e}tres entra\^{i}nables. Ce travail constitue la premi\`{e}re contribution de la th\`{e}se en \'{e}tablissant les fondations du routage intelligent \'{e}co\'{e}nerg\'{e}tique pour les r\'{e}seaux de capteurs sans fil (RCsF) de surveillance des catastrophes.}

\vspace{0.5cm}

\section{Abstract}

Wireless Sensor Networks (WSNs) deployed in disaster-monitoring scenarios require energy-efficient, low-latency multi-hop routing under strict resource constraints. Existing multi-agent reinforcement learning (MARL) approaches such as QMIX and QTRAN treat the routing problem independently, ignoring the network topology and suffering from memory explosion at scale. We propose \textbf{Graph-Attentive Policy Fusion (GAPF)}, a lightweight meta-learning framework that fuses frozen MARL expert policies through a Graph Convolutional Network (GCN) encoder and a Contextual Algorithm Selection (CAS) controller. GAPF encodes the sensor topology into a compact graph embedding at episode onset and selects the best-performing expert with only 1{,}473 trainable parameters. We evaluate GAPF across 7 scenarios spanning 20--100 sensors, 3 seeds each, totalling 21{,}000 test episodes. GAPF achieves near-perfect packet delivery ($\geq$98\%) across all scales, with +43\% to +111\% relative improvement over standalone MARL baselines. It consumes 2--3.5$\times$ less energy, maintains 3--9$\times$ better energy balance, and requires 13$\times$ less memory (951\,MB vs 13.2\,GB), passing all feasibility constraints. The scalability analysis reveals that GAPF's advantage grows superlinearly in sparse networks, where topology-aware relay coordination is most critical. Our results demonstrate that lightweight policy fusion with graph-structural awareness can replace heavyweight monolithic MARL for resource-constrained, safety-critical WSN routing.

\section{Introduction}

According to the European Environment Agency (EEA), a disaster is "a serious disruption of the functioning of society, causing widespread human, material or environmental losses, which exceed the ability of affected society to cope using only its own resources" \cite{noauthor_disaster_nodate}. This scourge is a real handicap for our society according to statistics. Indeed, 398 natural disasters were recorded in 2023  with 95,000 deaths and losses of around 380 billion US dollars \cite{burgueno_salas_global_nodate} \cite{noauthor_number_nodate}. Establishing a disaster monitoring policy thus appears to be essential. 

Wireless Sensor Networks (WSNs) are emerging as an indispensable resource for disaster monitoring \cite{al_qundus_wireless_2020}. They owe this to their ability to collect data more accurately with minimal human intervention \cite{bhatt_wireless_2021-1}. They are also resilient and maintain their connected state even if disasters have led to the destruction of infrastructure. In addition, they benefit from ease of scaling and configuration as required. In the case of wood fires, for example, detection methods such as watchtowers, satellite image processing methods, optical sensors, and methods based on digital cameras have disadvantages including inefficiency, consumption of energy, delay, accuracy, and implementation costs \cite{dampage_forest_2022-1}. This is where WSNs prove valuable, being self-configuring, independent of infrastructure, efficient, and low in energy consumption. Despite the undeniable advantages of WSNs, large-scale issues need to be considered, notably battery life due to the problem of energy consumption \cite{yarali_wireless_2020}. It is therefore important to reduce the waste of energy. This becomes even more critical as rising energy consumption leads to higher CO2 emissions, which is harmful to the environment \cite{thakare_study_2017}. 

Many algorithms have been proposed to solve the energy efficiency problem in this context. In sum, we are able to highlight some limitations in the reviewed literature:

\begin{itemize}
    \item A less realistic setting: In a disaster scenario, the wireless sensors static mobility consideration (\cite{priyanga_energy_2018}, \cite{mr_holistic_2023}) can fail since sensor position changes can be observed. Also, in a disaster setting, actions should be taken quickly to save lives and mitigate financial loss. For these reasons, real-time monitoring or at least near-real time monitoring should not be optional (\cite{maheshwari_energy_2021}, \cite{abualigah_swarm_2023}).
    \item The reward function modeling: Some factors significantly influence network energy efficiency and can’t be ignored. \cite{su_q-learning-based_2023} propose a routing algorithm based on Q-learning to decrease and balance the energy consumed by sensors. The reward function in this work is a sum of partial rewards. This can be misleading since the relationship between the partial rewards and the final one could be complex and so a simple weighted sum cannot capture all this complexity.
    \item The theoretical limitations: In its current formulation, QMIX is limited from a theoretical point of view (\cite{ding_packet_2022}, \cite{noauthor_multi-agent_nodate}). It is important to find the tradeoff between its good properties and some theoretical upgrade.
\end{itemize}
This work intends to overcome these limitations. 
\begin{enumerate}
    \item We propose \textbf{GAPF} (Graph-Attentive Policy Fusion), a lightweight
          meta-controller that sits on top of two independently pre-trained and
          frozen MADRL experts---QMIX and QTRAN---to capture the trade-off
          between a state-of-the-art value decomposition method and a
          theoretically more principled factorisation.
          GAPF encodes the WSN topology through a two-layer Graph Convolutional
          Network (GCN) and operates in two modes:
          \emph{CAS} (Contextual Algorithm Selection), which selects the
          best-suited expert for an entire episode based on the graph embedding,
          and \emph{Fusion}, which mixes the Q-value distributions of both
          experts at every time step via a learned softmax weighting.
          The meta-controller is trained end-to-end with REINFORCE using only
          $1{,}473$ parameters, making it extremely sample-efficient and
          deployable on resource-constrained devices.

    \item We develop a bi-level reward structure built on a
          \textbf{Monotonic Attention Network}.  At the inner level, each
          sensor receives a four-dimensional reward vector capturing
          (i)~angular progress toward the base station,
          (ii)~transmission energy consumption,
          (iii)~dispersion of remaining energy across the network, and
          (iv)~buffer occupancy.
          A query--key attention mechanism learns objective-specific weights
          $\alpha_i = \operatorname{softmax}(\mathbf{q} \cdot \mathbf{k}_i
          / \sqrt{d})$, and a monotonic multilayer network
          ($4 \!\to\! 16 \!\to\! 16 \!\to\! 1$, all weights $\geq 0$)
          scalarises the weighted vector into a single reward, guaranteeing
          that component-wise improvement always increases the scalar signal.
          At the outer level, the attention parameters are updated after each
          episode via \textbf{Evolution Strategies} to maximise a meta-objective
          $J$ that combines energy efficiency and energy balance, decoupling
          reward shaping from the RL policy gradient.

    \item We design a custom gym environment that is more realistic for a
          disaster scenario because it considers the dynamic mobility of nodes
          and factors which imbalance and increase the overall energy consumption
          of the network.
\end{enumerate}



\section{Background and Related Work}
In recent years, a lot of research has been done on the routing problem. Moussa et al. \cite{moussa_reinforcement_2023} propose a Reinforcement Learning (RL) based routing protocol for Software Defined Network (SDN)-enabled WSN forest fire detection. The first step of the protocol was to design a clustering algorithm that introduces a delay in the re-clustering process to reduce energy. Secondly, an optimal path is defined by the SDN controller with the help of RL. Their work shows good results in network performance, however, it suffers from some drawbacks. Mainly, there is a lack of generalization in the RL model and a lack of being able to capture an important part in the factors which are involved in the reward construction. 

Mishra et al. \cite{mishra_meta-heuristic-based_2021} propose the Energy-Balanced Grey Wolf Optimizer (EBGWO) for the clustering routing issue in Software-Defined WSN (SD-WSN). In their work, the selection of control nodes is tackled with the Gray-Wolf Optimization approach for energy balance. EBGWO succeeded in extending the network lifespan. The limitation however is the static consideration of the network \cite{isyaku_software_2024}. 

Biabani et al. \cite{biabani_energy-efficient_2020} provide a clustering and routing approach able to manage the nodes’ energy consumption knowing the disaster environment characteristics. Their model is presented with hybrid Particle Swarm Optimization (PSO) and Harmony Search Algorithm to select the cluster head. Their multi-hop routing method is also based on PSO. Therefore, in disasters settings, the model is more efficient than the state-of-the-art approaches considering residual energy and other parameters. However, there are important delays regarding routing decisions \cite{noauthor_survey_nodate} and the nodes are static in the network \cite{biabani_energy-efficient_2020} which can be an issue in a disaster context. 

Hosseinzadeh et al. \cite{hosseinzadeh_hybrid_2022} introduce a swarm intelligence-based clustering method to have cluster heads distributed more evenly. To select the cluster head and intermediate nodes for the routing process, authors use a firefly optimization approach and an aquila optimizer algorithm. This method improves the energy of the system. Nevertheless, the dynamic relocation of nodes needs to be explored \cite{senkumar_enhanced_2024}. 

Dehkordi et al. \cite{ghorbani_dehkordi_cluster_2023} propose a cluster-based routing model in WSN by using mobile sinks to decrease the nodes energy consumption and improve the network longevity. When comparing their method with other models, it comes at its superiority in terms of energy consumption among others. But the fact that mobile sinks are involved in their scheme, this may lead to some delays \cite{senkumar_enhanced_2024}. 

Gopalan et al. \cite{harihara_gopalan_energy_2024} elaborate a routing protocol where the Bacterial Foraging Optimization with Harmony Search Algorithm (BFO-HSA) is used for sensor nodes classification. The Cross-Layer-Based Opportunistic Routing Protocol (CORP) is considered to find the quickest path. Although the results demonstrate that their model outperforms the state of the art on certain quality of service parameters, it must also be recognized, as the authors point out, that the focus was on static sensor nodes \cite{harihara_gopalan_energy_2024}. 

Kodati et al. \cite{kodati_hybrid_2024} propose a Hybrid Grasshopper and Harris Hawk Optimization Algorithm-based energy efficient routing protocol for optimal selection of clusters heads. At this effect, they consider a fitness function which considers many factors. Outcomes of this approach are better in terms of residual energy and throughput compared to the baseline methods. Yet, with this approach, there may be an energy imbalance \cite{kodati_hybrid_2024}.

Wang et al. \cite{wang_elite_2020} build an elite hybrid metaheuristic optimization algorithm-based routing algorithm to maximize the network lifetime of the WSN with a sink node. They successfully improve the survival time of the WSN in comparison to the state-of-the-art. However, the process of Cluster Head selection needs room for improvement \cite{kodati_hybrid_2024}. 

Chaurasia et al. develop \cite{chaurasia_mocraw_2023} a Meta-Heuristic Optimized Cluster Head Selection-Based Routing Algorithm for WSNs to minimize the energy consumption of sensors and increase the speed of the transmission of data. Underneath their model is a Dragonfly Algorithm for CH and path selection. It outperforms other models regarding the energy efficiency. However, the delay is higher \cite{wilson_optimizing_2024}. 

Yao et al. \cite{yao_energy-efficient_2022} design an Energy-Efficient Routing Protocol based on Multi-Threshold Segmentation to improve the clustering process (cluster creation and the selection of CH). However, the CH selection is not properly done due to the problem of the overhead and clustering \cite{rao_optimized_2024}. 

Manoharan et al. \cite{manoharan_optimal_2024} propose an energy efficient routing in WSN using adaptive entropy bald eagle search optimization and density based adaptive soft clustering to prove the optimal data transfer in the network. The author’s approach reaches better performance than other methods. A drawback is the non-consideration of mobile nodes \cite{manoharan_optimal_2024}. 

Rai et al. \cite{rai_feec_2023} suggest a fuzzy Super Cluster Head (SCH) selection approach by considering four factors to reduce the dissipation of energy and improve the network longevity. Although it is true that this method outperforms some protocols, it should also be noted that under this approach there is a hot spot issue \cite{dinesh_hbo-sroa_2024}. 

Prakash et al. \cite{prakash2025enhanced} introduce an energy-efficient, cluster-based routing protocol — denoted SHO-CH — which leverages a Spotted Hyena Optimization strategy for cluster-head selection, and is tailored for heterogeneous wireless sensor networks (WSNs). SHO-CH optimize cluster formation and routing in order to reduce energy consumption across sensor nodes, thereby extending the operational lifetime of the network and enhancing energy-efficiency. The evaluation of the proposed model demonstrates improvements in stability period, overall network lifetime, and throughput compared with baseline protocols. However, the assumption of static sensor nodes and reliance on single-hop intra-cluster communication — as well as the computational overhead introduced by the iterative optimization process — limit the scheme’s realism and applicability, especially in dynamic and unpredictable environments such as disaster-scenarios. 

Ataoua et al. \cite{ataoua_chernobyl_2025} introduce a novel clustering protocol for wireless sensor networks (WSNs), leveraging the Chernobyl Disaster Optimizer (CDO) metaheuristic to enhance energy efficiency and extend network lifetime. Their approach focuses on optimizing cluster-head selection and communication paths by minimizing intra-cluster distances and balancing energy consumption across nodes. Evaluated through simulations against standard protocols such as LEACH, the proposed method exhibits marked improvements in residual energy retention, reduced node failure rates, and prolonged network operation. Nonetheless, the study is constrained by its reliance on static simulation environments, omitting evaluation under disaster-specific dynamics such as node mobility, cascading failures, and communication interruptions. Additionally, the protocol overlooks key performance metrics including latency, packet delivery ratio, and security—factors essential to robust deployment in disaster monitoring contexts.

Qamar et al. \cite{qamar_hybrid_2025} propose a hybrid routing framework that combines deterministic clustering with neural network–based energy prediction and energy-aware path selection to improve the efficiency and operational lifetime of wireless sensor networks (WSNs). The methodology integrates a static clustering mechanism derived from LEACH-D with an artificial neural network trained to estimate node energy levels, enabling routing decisions that prevent excessive load on low-energy nodes and promote balanced energy consumption across the network. Simulation results show reductions in energy consumption and notable improvements in network lifetime and packet delivery ratio compared with conventional protocols such as LEACH and related approaches. However, the framework assumes static network topologies and does not account for node mobility or failure dynamics commonly encountered in disaster scenarios. Furthermore, the neural network model is trained offline, which limits its adaptability to real-time environmental variations, and the absence of integrated security or trust mechanisms restricts its applicability in adversarial or disaster-prone deployments.

Rahman et al. \cite{rahman_distributed_2025} propose a blockchain-enabled architecture for wireless sensor networks (WSNs), aimed at detecting malicious nodes and preserving data integrity in critical scenarios such as disaster monitoring. The framework integrates a hybrid consensus model that combines smart contracts with a Proof of Information (PoI) mechanism to validate sensor data and isolate compromised nodes. Simulations report a data validity rate of 97.37\%, highlighting the benefits of separating monitoring and control sensors to improve real-time responsiveness. However, the design does not incorporate energy-aware routing or clustering strategies, and its scalability in large-scale deployments remains unaddressed. Moreover, the protocol lacks evaluation under high-mobility or failure-prone conditions typical of disaster environments, limiting its practical applicability.

Rao et al. \cite{rao_optimized_2024} propose a secure, cluster-based routing protocol for IoT-enabled wireless sensor networks (WSNs), incorporating a lightweight blockchain framework to ensure data integrity and resist tampering in decentralized deployments. The protocol integrates energy-aware clustering with hash-based data validation and smart contract mechanisms, aiming to enhance trustworthiness without compromising efficiency. Experimental evaluations report improved energy consumption, lower packet loss, and heightened data security relative to conventional routing approaches. However, the applicability of the proposed system to disaster monitoring remains limited, as the study is confined to agricultural use cases and does not address critical challenges such as node mobility, failure resilience, or intermittent connectivity. Furthermore, the added communication overhead introduced by the blockchain layer may hinder latency and scalability, especially in time-sensitive emergency response scenarios.

Sefati et al. \cite{sefati_federated_2025} propose AFRL-HGS, a federated reinforcement learning framework combined with hybrid optimization techniques to support secure and adaptive routing in wireless sensor networks (WSNs). In this approach, sensor nodes act as reinforcement learning agents that collaboratively train routing policies through federated aggregation without sharing raw data, thereby preserving privacy and reducing communication overhead, while the hybrid optimization component accelerates convergence and stabilizes routing decisions. Simulation-based evaluation demonstrates improvements in packet delivery ratio, energy efficiency, and resilience to node failures compared with centralized reinforcement learning and conventional routing protocols. However, the framework assumes periodic synchronization and relatively stable connectivity for aggregating local models—conditions that may not hold in disaster environments characterized by network partitioning and high mobility—and it does not integrate additional trust or blockchain-based safeguards, potentially leaving the collaborative learning process vulnerable to model poisoning or adversarial manipulation.

Table~\ref{art1:tab:comparative_metrics} shows a comparative Analysis of Selected Routing Algorithms (2024–2025).

\begin{center}
\captionof{table}{Comparative Analysis of Selected Routing Algorithms (2024–2025)}
\resizebox{\textwidth}{!}{
\begin{tabular}{|p{4cm}|p{3.2cm}|c|c|c|c|}
\hline
\textbf{Article (Author, Year)} & \textbf{Approach} & \textbf{Energy Efficiency} & \textbf{PDR} & \textbf{Latency} & \textbf{Mobility Support} \\
\hline
Ataoua et al. (2025) & Chernobyl Disaster Optimizer (CDO)-based clustering & High & Not evaluated & Not analyzed & No \\
\hline
Rahman et al. (2025) & Blockchain-secure routing with PoI consensus & Not addressed & High (97.37\% data validity) & Moderate (blockchain overhead not quantified) & No \\
\hline
Qamar \& Munir (2025) & Neural-assisted clustering with energy-aware routing & High & Improved & Not quantified & No \\
\hline
Rao et al. (2024) & Secure cluster-based routing with blockchain integrity checking & Moderate & Improved & Slightly increased due to blockchain & Limited \\
\hline
Sefati et al. (2025) & Federated RL with hybrid optimization (AFRL-HGS) & High & High & Low (fast convergence) & Partial (assumes stable sync) \\
\hline
\end{tabular}
}
\label{art1:tab:comparative_metrics}
\end{center}


This literature review shows the extensive work that has been done to achieve the energy efficiency goal of routing protocols and the limitations that they have faced which motivates our current work.


\section{Methodology}
\subsection{Problem description}
The basis of the proposed hybrid model consists of two Centralized Training and Decentralized Execution (CTDE) models named QMIX and QTRAN. As such, a formalization is needed in a context of MADRL. We assume there is a Base Station (BS) and \textit{N} sensors nodes with the following characteristics:

\begin{equation} 
S_k \equiv 
\begin{cases}
CoverageRadius_k \text{ } \left(Rad_k\right)  \\
Position_k \\
RemainingEnergy_k \text{ } \left(RE_k\right) \\
ConsumptionEnergy_k \text{ } \left(CE_k\right) \\
Number of packets \text{ } \left(N_k\right) 
\end{cases}
\end{equation}

with $Rad_k$, $Position_k$, $RE_k$, $CE_k$, $N_k$ being respectively the coverage radius, the coordinates in the environment to sense, the remaining energy, the consumption energy and the number of packets held by the $k^{th}$ sensor.  

Each sensor sends data to the base station through a multi hop process.

We assume that all sensors have the same coverage radius denoted $R$ and at the beginning of each episode of the training, they have one packet to transmit. We set $E_{max}$ to be the initial energy for all the sensors. 

The observation space is defined as follows:
\begin{itemize}
    \item Remaining energy for all sensors.
    \item Consumption energy for all sensors.
    \item Positions in the environment for all sensors: A two-dimensional space is considered.
    \item Number of packets for all sensors.
\end{itemize}	
    
The routing objective is threefold:
\begin{enumerate}[nosep]
    \item \textbf{Maximize packet delivery ratio} --- deliver as many sensor readings as possible to the base station, the primary goal of any data-gathering WSN;
    \item \textbf{Minimize total energy consumption} --- reduce the cumulative transmission and reception energy across all sensors to extend network operational lifetime;
    \item \textbf{Balance energy usage across sensors} --- minimize the standard deviation of remaining energy, preventing premature death of overloaded relay nodes.
\end{enumerate}
These objectives are inherently conflicting: delivering more packets requires more transmissions, which increases energy consumption. The proposed framework addresses this tension through a bi-level optimization architecture. The \emph{inner loop} (MARL policy) receives a scalar reward that combines all three objectives via four individual reward components --- relay angle, energy consumption, energy dispersion, and receiver load --- weighted by a Monotonic Q$\cdot$K Attention Network. The \emph{outer loop} (Evolution Strategies) tunes the attention network's parameters to maximize a meta-objective $J$ that focuses on energy efficiency and energy balance, since packet delivery is already strongly incentivized by a dominant delivery bonus ($+100$) in the inner-loop reward. This separation ensures that the RL agents prioritize delivery during training while the scalarizer progressively learns to route energy-efficiently.

\subsection{The Energy model}
From \cite{su_q-learning-based_2023} and \cite{guo_optimizing_2019}, the energy model is defined as:
\begin{equation} 
\left\{
\begin{aligned}
  E_T\left(M, d\right) &= M \times a \times \left(E_{elec} + E_{amp} \times d^2\right) \\
  E_R\left(M\right) &= M \times a \times E_{elec}
\end{aligned}
\right.
\end{equation}

where $E_{elec} = 50\frac{nJ}{bit}$, $E_{amp} = 100\frac{pJ}{bit \times m^2}$, $E_T\left(M, d\right)$ is the energy that a sensor consumes to transmit $M$ packets (each packet contains $a$ amount of information) to a receiver at the distance $d$, $E_R\left(M\right)$  is the energy that a sensor consumes to receive $M$ packets from the sender. 

\subsection{Reward Scalarization via Monotonic Q$\cdot$K Attention Network}
\label{art1:sec:reward_scalarization}

In cooperative MARL for WSN routing, agents must simultaneously optimize energy efficiency, energy balance, directional progress toward the base station, and load distribution. Fixed linear scalarizations require manual weight tuning that does not generalize across topologies. We introduce a \emph{Monotonic Q$\cdot$K Attention Network} that \textbf{(i)} learns objective importance weights via scaled dot-product attention~\cite{vaswani_attention_2023}, \textbf{(ii)} guarantees monotonicity so that Pareto-improving actions are never penalized~\cite{sill_monotonic_1997}, and \textbf{(iii)} adapts online through an Evolution Strategies (ES) outer loop~\cite{salimans2017evolution}.

\subsubsection{Individual Reward Objectives}
\label{art1:sec:individual_rewards}

At each step $t$, when sensor $i$ relays packets to neighbor $j$, the environment computes a four-dimensional reward vector $\mathbf{r}_i^{(t)} = [r_1, r_2, r_3, r_4] \in [0,1]^4$:
\begin{itemize}[nosep]
    \item \textbf{Directional progress} ($r_1$): $r_1 = 1 - |\theta_{ij}|/\pi$, where $\theta_{ij}$ is the angle between the relay direction $\overrightarrow{ij}$ and the ideal direction toward the base station (BS) $\overrightarrow{i\,\mathrm{BS}}$. Perfect alignment yields $r_1 = 1$; opposite direction yields $r_1 = 0$.
    \item \textbf{Energy efficiency} ($r_2$): $r_2 = 1 - (E_{\mathrm{tx}}(p_i, d_{ij}) + E_{\mathrm{rx}}(p_i))/(E_{\mathrm{tx}}^{\max} + E_{\mathrm{rx}}^{\max})$, where $E_{\mathrm{tx}}(p_i, d_{ij})$ and $E_{\mathrm{rx}}(p_i)$ are the transmission and reception energies for $p_i$ packets over distance $d_{ij}$, normalized by worst-case values $E_{\mathrm{tx}}^{\max} = E_{\mathrm{tx}}(N \cdot p_0, R)$ and $E_{\mathrm{rx}}^{\max} = E_{\mathrm{rx}}(N \cdot p_0)$, with $N$ the number of sensors, $p_0$ the initial packet count per sensor, and $R$ the communication radius.
    \item \textbf{Energy balance} ($r_3$): $r_3 = 1 - \sigma(\mathbf{e}_{\mathrm{rem}})/(E_0/2)$, where $\sigma(\mathbf{e}_{\mathrm{rem}})$ is the standard deviation of remaining energies across all sensors, $E_0$ is the initial energy budget per sensor, and $E_0/2$ is the theoretical maximum standard deviation (when half the sensors are fully charged and half depleted).
    \item \textbf{Load balancing} ($r_4$): $r_4 = 1 - p_j/(N \cdot p_0)$, where $p_j$ is the current packet count at receiver node $j$.
\end{itemize}
All four components are clipped to $[0,1]$.

\subsubsection{Monotonic Q$\cdot$K Attention Architecture}
\label{art1:sec:monotonic_architecture}

The scalarization function $f_\phi: [0,1]^4 \to \mathbb{R}$ maps the reward vector $\mathbf{r}$ to a scalar via two stages.

\paragraph{Stage~1: Attention weighting.}
A learnable query vector $\mathbf{q} \in \mathbb{R}^d$ and a key matrix $\mathbf{K} = [\mathbf{k}_1, \dots, \mathbf{k}_4]^\top \in \mathbb{R}^{4 \times d}$ (with embedding dimension $d = 16$) produce attention weights $\alpha_i = \mathrm{softmax}(\mathbf{K}\mathbf{q}/\sqrt{d})_i$, yielding the weighted input $\mathbf{r}' = \boldsymbol{\alpha} \odot \mathbf{r}$ (element-wise product). Because $\boldsymbol{\alpha}$ depends only on the learnable parameters $(\mathbf{q}, \mathbf{K})$ and not on $\mathbf{r}$, and since $\alpha_i > 0$ for all $i$ (due to softmax), the weighting preserves the component-wise partial order: $\mathbf{r}^{(1)} \leq \mathbf{r}^{(2)}$ implies $\mathbf{r}'^{(1)} \leq \mathbf{r}'^{(2)}$.

\paragraph{Stage~2: Monotonic network.}
The weighted input $\mathbf{r}'$ passes through a three-layer feed-forward network ($4 \to 16 \to 16 \to 1$) with softplus activations $\sigma(x) = \log(1 + e^x)$ and non-negative weight matrices enforced via $\mathbf{W}_\ell = \mathrm{softplus}(\mathbf{W}_\ell^{\mathrm{raw}}) \geq 0$. Since the softplus derivative $\sigma'(x) = (1+e^{-x})^{-1} \in (0,1)$ is strictly positive and all weights are non-negative, the Jacobian satisfies $\partial f / \partial r_i \geq 0$ for every component. By the chain rule over all layers, the full network $f_\phi = g_\psi \circ (\boldsymbol{\alpha} \odot \cdot)$ is monotonically non-decreasing: any Pareto-dominating action always receives a higher scalar reward. The output is centered as $\hat{f}_\phi(\mathbf{r}) = f_\phi(\mathbf{r}) - f_\phi(\mathbf{0})$ to improve signal-to-noise ratio, where the baseline $f_\phi(\mathbf{0})$ is recomputed once per episode.

\subsubsection{Composite Step Reward}
\label{art1:sec:composite_reward}

The final reward for sensor $i$ at step $t$ depends on its action: \textbf{(a)}~delivery to the BS yields $r_i^{(t)} = \beta_{\mathrm{del}} \cdot p_i$, where $\beta_{\mathrm{del}} = 100$ is a delivery bonus and $p_i$ the number of packets delivered (bypassing scalarization, as delivery is the terminal objective); \textbf{(b)}~relay to neighbor $j \neq \mathrm{BS}$ yields:
\begin{equation}
    r_i^{(t)} = \underbrace{\hat{f}_\phi(\mathbf{r}_i^{(t)}) \cdot \lambda_{\mathrm{attn}}}_{\text{attention-scalarized}} \;+\; \underbrace{\frac{d(i, \mathrm{BS}) - d(j, \mathrm{BS})}{R} \cdot 0.5}_{\text{progress shaping}}
    \label{art1:eq:relay_reward}
\end{equation}
where $\lambda_{\mathrm{attn}} = 0.01$ scales the attention output to prevent relay-loop farming, and $d(\cdot, \mathrm{BS})$ denotes the Euclidean distance to the BS; \textbf{(c)}~invalid actions receive penalties of $-0.1$ (self-loops, out-of-range) to $-0.5$ (insufficient energy); inactive sensors receive $0$.

\subsubsection{Evolution Strategies Meta-Learning}
\label{art1:sec:es_learning}

The attention network parameters $\phi = (\mathbf{q}, \mathbf{K}, \{\mathbf{W}_\ell^{\mathrm{raw}}, \mathbf{b}_\ell\}_\ell)$ are trained via an outer-loop ES~\cite{salimans2017evolution} operating at episode granularity, without differentiating through the RL episode. At episode end, a meta-objective $J \in [0,2]$ evaluates the true network-level goals:
\begin{equation}
    J = \underbrace{\left(1 - \frac{\sum_{i=1}^{N}(E_0 - e_i^{\mathrm{rem}})}{N \cdot E_0}\right)}_{\text{energy efficiency}} + \underbrace{\left(1 - \frac{\sigma(\mathbf{e}_{\mathrm{rem}})}{E_0/2}\right)}_{\text{energy balance}}
    \label{art1:eq:meta_objective}
\end{equation}
where $e_i^{\mathrm{rem}}$ is the remaining energy of sensor $i$ at episode end. Each episode: \textbf{(1)}~parameters are perturbed as $\phi' = \phi + \sigma_{\mathrm{ES}}\,\boldsymbol{\varepsilon}$ with $\boldsymbol{\varepsilon} \sim \mathcal{N}(\mathbf{0}, \mathbf{I})$ and noise scale $\sigma_{\mathrm{ES}} = 0.01$; \textbf{(2)}~the episode executes with $\phi'$ and $J(\phi')$ is computed; \textbf{(3)}~the advantage $A = J(\phi') - \bar{J}$ is evaluated against an exponential moving average baseline $\bar{J}$ (decay $\beta = 0.9$), and parameters are updated as $\phi \leftarrow \phi + \eta_{\mathrm{ES}} \cdot A \cdot \boldsymbol{\varepsilon}$ with learning rate $\eta_{\mathrm{ES}} = 0.01$ (skipped when $|A| < 0.01$). This bi-level optimization decouples inner-loop RL training from outer-loop reward shaping. Key embeddings for energy ($\mathbf{k}_2$) and dispersion ($\mathbf{k}_3$) are initialized close to $\mathbf{q}$ (higher initial attention), encoding the prior that energy-related objectives are most important, while the ES procedure discovers the optimal weighting. The learned weights $\boldsymbol{\alpha}$ provide direct interpretability of objective importance at any point during training.




\subsection{GAPF: Graph-Attentive Policy Fusion}
\label{art1:sec:gapf}

We introduce \textbf{GAPF} (Graph-Attentive Policy Fusion), a lightweight meta-algorithm that selects, at the episode level, the best-performing MARL expert for a given WSN topology. GAPF uses a graph neural network (GNN) to encode the spatial structure of the sensor network and a meta-controller trained via REINFORCE~\cite{williams1992simple} to choose between two pre-trained, frozen experts (QMIX and QTRAN).

\subsubsection{Problem Formulation}
\label{art1:sec:gapf_formulation}

Let $\mathcal{E} = \{\pi_{\text{QMIX}}, \pi_{\text{QTRAN}}\}$ denote pre-trained, frozen MARL policies for the WSN routing Dec-POMDP. At the start of each episode~$k$, the WSN defines a communication graph $\mathcal{G}_k = (\mathcal{V}, \mathcal{E}_k)$, where nodes $\mathcal{V} = \{1, \dots, N\}$ are sensors and an edge $(i,j) \in \mathcal{E}_k$ exists iff $\|p_i - p_j\|_2 \leq R$, with $p_i \in \mathbb{R}^2$ the position of sensor~$i$ and $R$ the communication radius. GAPF learns a parameterised policy $\mu_\theta$ mapping the topology $\mathcal{G}_k$ to a selection $a^{\text{meta}}_k \in \{\text{QMIX}, \text{QTRAN}\}$, maximising the expected episodic return:
\begin{equation}
    \max_\theta \; \mathbb{E}_{\mathcal{G}_k}\!\left[\, G_k \,\middle|\, \pi_{a^{\text{meta}}_k}\right], \quad a^{\text{meta}}_k \sim \mu_\theta(\cdot \mid \mathcal{G}_k).
    \label{art1:eq:gapf_objective}
\end{equation}

\subsubsection{Architecture}
\label{art1:sec:gapf_architecture}

The GAPF model $\mu_\theta$ consists of a \emph{GCN encoder} producing a fixed-dimensional graph embedding, and a \emph{CAS controller} (Contextual Algorithm Selection) mapping this embedding to a Bernoulli selection probability (Figure~\ref{art1:fig:gapf_architecture}).

\paragraph{Graph construction.}
Given the $N$ sensor positions $\{p_i\}_{i=1}^N$ and communication radius~$R$, the binary adjacency matrix $\mathbf{A} \in \{0,1\}^{N \times N}$ is defined as $A_{ij} = \mathbf{1}[\|p_i - p_j\| \leq R] \cdot \mathbf{1}[i \neq j]$, then symmetrically normalised following~\cite{kipf2017semi}:
\begin{equation}
    \hat{\mathbf{A}} = \tilde{\mathbf{D}}^{-1/2}\,\tilde{\mathbf{A}}\,\tilde{\mathbf{D}}^{-1/2}, \quad \tilde{\mathbf{A}} = \mathbf{A} + \mathbf{I}_N, \quad \tilde{D}_{ii} = \textstyle\sum_j \tilde{A}_{ij}.
    \label{art1:eq:norm_adj}
\end{equation}
Each sensor $i$ is described by a feature vector $\mathbf{x}_i \in \mathbb{R}^4$:
\begin{equation}
    \mathbf{x}_i = \left[\; \frac{p_i^{(x)}}{L},\;\; \frac{p_i^{(y)}}{L},\;\; \frac{\|p_i - p_{\text{BS}}\|}{\max_j \|p_j - p_{\text{BS}}\|},\;\; \frac{\deg(i)}{\max_j \deg(j)} \;\right]^\top,
    \label{art1:eq:node_features}
\end{equation}
where $L$ is the field side length, $p_{\text{BS}}$ is the base station position, and $\deg(i) = \sum_j A_{ij}$ is the node degree. All features are normalised to $[0, 1]$.

\paragraph{GCN encoder.}
The node feature matrix $\mathbf{X} = [\mathbf{x}_1, \dots, \mathbf{x}_N]^\top \in \mathbb{R}^{N \times 4}$ is processed by a two-layer GCN~\cite{kipf2017semi}:
\begin{equation}
    \mathbf{H}^{(1)} = \text{ReLU}(\hat{\mathbf{A}}\,\mathbf{X}\,\mathbf{W}_1), \quad
    \mathbf{Z} = \hat{\mathbf{A}}\,\mathbf{H}^{(1)}\,\mathbf{W}_2, \quad
    \mathbf{z} = \tfrac{1}{N}\textstyle\sum_{i=1}^{N} \mathbf{Z}_{i,:} \in \mathbb{R}^{d}
    \label{art1:eq:gcn}
\end{equation}
where $\mathbf{W}_1 \in \mathbb{R}^{4 \times d}$, $\mathbf{W}_2 \in \mathbb{R}^{d \times d}$, and $d = 16$ is the embedding dimension. The mean-pool readout is permutation-invariant and handles variable numbers of sensors without architectural changes.

\paragraph{CAS controller.}
The graph embedding $\mathbf{z}$ is fed to a two-layer MLP that outputs the probability of selecting QMIX:
\begin{equation}
    p_{\text{QMIX}} = \sigma\!\left(\mathbf{w}_2^\top\,\text{ReLU}\!\left(\mathbf{W}_c\,\mathbf{z} + \mathbf{b}_c\right) + b_2\right),
    \label{art1:eq:cas_controller}
\end{equation}
where $\mathbf{W}_c \in \mathbb{R}^{h \times d}$, $\mathbf{b}_c \in \mathbb{R}^h$, $\mathbf{w}_2 \in \mathbb{R}^h$, $b_2 \in \mathbb{R}$, $h = 64$ is the hidden dimension, and $\sigma$ is the sigmoid function. During training, $a^{\text{meta}} \sim \text{Bernoulli}(p_{\text{QMIX}})$; at evaluation, $a^{\text{meta}} = \text{QMIX}$ if $p_{\text{QMIX}} > 0.5$, QTRAN otherwise. The total number of trainable parameters is:
\begin{equation}
    |\theta| = \underbrace{4 \times 16 + 16 \times 16}_{\text{GCN: } 320} \;+\; \underbrace{16 \times 64 + 64 + 64 \times 1 + 1}_{\text{CAS: } 1{,}153} = 1{,}473.
    \label{art1:eq:param_count}
\end{equation}
This is approximately $26\times$ smaller than a single RNN agent ($38{,}983$ parameters), making GAPF negligible in computational overhead.

% --- FIGURE (tikzpicture unchanged) ---
\begin{center}
\resizebox{\linewidth}{!}{%
\begin{tikzpicture}[
    node distance=0.9cm,
    box/.style={draw, rounded corners, minimum height=0.8cm, minimum width=4.2cm, align=center, font=\small},
    data/.style={draw, rounded corners, minimum height=0.6cm, minimum width=3.0cm, align=center, fill=gray!10, font=\small},
    arrow/.style={-{Stealth[length=2.5mm]}, thick},
    dasharrow/.style={-{Stealth[length=2.5mm]}, thick, dashed},
    label/.style={font=\scriptsize, text=black!60},
]
\node[data] (topo) {WSN Topology $\{p_i\}_{i=1}^N,\; R$};
\node[box, below=of topo, fill=blue!8] (graph) {Graph Construction\\{\scriptsize $\hat{\mathbf{A}} \in \mathbb{R}^{N \times N}$, $\mathbf{X} \in \mathbb{R}^{N \times 4}$}};
\node[box, below=of graph, fill=green!12] (gcn) {GCN Encoder {\scriptsize (2-layer, $d\!=\!16$)}};
\node[box, below=of gcn, fill=green!12] (pool) {Mean-Pool Readout};
\node[data, below=of pool] (embed) {$\mathbf{z} \in \mathbb{R}^{16}$};
\node[box, below=of embed, fill=orange!15] (cas) {CAS Controller {\scriptsize (MLP: $16 \!\to\! 64 \!\to\! 1$)}};
\node[data, below=of cas] (prob) {$p_{\text{QMIX}} = \sigma(\text{logit})$};
\node[box, below=of prob, fill=red!10] (decision) {Episode-level Selection};
\node[box, below left=1.0cm and 0.4cm of decision, fill=cyan!12, minimum width=2.0cm] (qmix) {$\pi_{\text{QMIX}}$\\{\scriptsize frozen}};
\node[box, below right=1.0cm and 0.4cm of decision, fill=cyan!12, minimum width=2.0cm] (qtran) {$\pi_{\text{QTRAN}}$\\{\scriptsize frozen}};
\node[box, below=2.4cm of decision, fill=yellow!12, minimum width=4.8cm] (env) {WSN Environment\\{\scriptsize 100 steps $\to$ return $G_k$}};
\draw[arrow] (topo) -- (graph);
\draw[arrow] (graph) -- (gcn);
\draw[arrow] (gcn) -- node[right, label] {$\mathbf{Z} \in \mathbb{R}^{N \times 16}$} (pool);
\draw[arrow] (pool) -- (embed);
\draw[arrow] (embed) -- (cas);
\draw[arrow] (cas) -- (prob);
\draw[arrow] (prob) -- (decision);
\draw[dasharrow] (decision) -- node[left, label] {$p > 0.5$} (qmix);
\draw[dasharrow] (decision) -- node[right, label] {$p \leq 0.5$} (qtran);
\draw[arrow] (qmix) |- (env);
\draw[arrow] (qtran) |- (env);
\draw[arrow, red!70!black, thick]
    (env.west) -- ++(-1.5, 0) |- (cas.west)
    node[pos=0.25, left, font=\scriptsize, text=red!70!black, align=right]
    {REINFORCE\\$\nabla_\theta J = (G_k - b_k)\,\nabla_\theta \log \mu_\theta$};
\draw[decorate, decoration={brace, amplitude=4pt, mirror}, thick, black!50]
    let \p1=(gcn.north east), \p2=(prob.south east) in
    (\x1+0.8cm, \y1) -- (\x1+0.8cm, \y2)
    node[midway, right=8pt, font=\scriptsize, text=black!60, align=left] 
    {$|\theta| = 1{,}473$\\trainable};
\end{tikzpicture}
}
\captionof{figure}{GAPF architecture (CAS mode). Sensor positions and communication radius define a graph encoded by a two-layer GCN with mean-pool readout ($\mathbf{z} \in \mathbb{R}^{16}$). The CAS controller maps $\mathbf{z}$ to a selection probability $p_{\text{QMIX}}$. The chosen frozen expert executes the full episode. The return $G_k$ feeds back (red arrow) as a REINFORCE signal to update only the $1{,}473$ meta-controller parameters.}
\label{art1:fig:gapf_architecture}
\end{center}

\subsubsection{Training Procedure}
\label{art1:sec:gapf_training}

QMIX and QTRAN are trained independently; their parameters are then frozen. Only the GAPF meta-controller ($1{,}473$ parameters) is trained via REINFORCE~\cite{williams1992simple} with an EMA baseline. Let $G_k$ denote the episodic return from executing the selected expert $\pi_{a^{\text{meta}}_k}$ in episode~$k$. The policy gradient is:
\begin{equation}
    \nabla_\theta J(\theta) = \mathbb{E}\!\left[\left(G_k - b_k\right) \nabla_\theta \log \mu_\theta\!\left(a^{\text{meta}}_k \mid \mathcal{G}_k\right)\right],
    \label{art1:eq:reinforce}
\end{equation}
where $b_k$ is the EMA baseline updated as $b_{k+1} = 0.95\, b_k + 0.05\, G_k$. Since the selection is Bernoulli, $\log \mu_\theta = \log p_{\text{QMIX}}$ if QMIX is chosen and $\log(1 - p_{\text{QMIX}})$ otherwise. Parameters are updated via Adam with learning rate $\alpha = 3 \times 10^{-3}$ and gradient clipping $\|\nabla_\theta\|_2 \leq 1.0$. The complete procedure is summarised in Algorithm~\ref{art1:alg:gapf_training}.

\begin{algorithm}[H]
\caption{GAPF Training (CAS Mode)}
\label{art1:alg:gapf_training}
\begin{algorithmic}[1]
\REQUIRE Frozen experts $\pi_{\text{QMIX}}$, $\pi_{\text{QTRAN}}$; WSN environment $\mathcal{M}$; training episodes $K$
\ENSURE Trained meta-controller parameters $\theta$
\STATE Initialise $\theta$ (Xavier uniform); Adam optimiser; $b \leftarrow \text{None}$
\FOR{$k = 1, \dots, K$}
    \STATE Reset $\mathcal{M}$; build $\hat{\mathbf{A}}_k$ (Eq.~\ref{art1:eq:norm_adj}), $\mathbf{X}_k$ (Eq.~\ref{art1:eq:node_features}); compute $\mathbf{z}_k$ via GCN (Eq.~\ref{art1:eq:gcn})
    \STATE $p_{\text{QMIX}} \leftarrow \text{CAS}(\mathbf{z}_k)$; \; sample $a^{\text{meta}}_k \sim \text{Bernoulli}(p_{\text{QMIX}})$
    \STATE Execute $\pi_{a^{\text{meta}}_k}$ for full episode; collect return $G_k$
    \STATE $b \leftarrow G_k$ \textbf{if} $b = \text{None}$; \; $\delta_k \leftarrow G_k - b$; \; $b \leftarrow 0.95\,b + 0.05\,G_k$
    \IF{$|\delta_k| > 10^{-8}$}
        \STATE $\theta \leftarrow \theta - \alpha \cdot \text{Adam}(\text{clip}(\nabla_\theta[-\delta_k \cdot \log \mu_\theta(a^{\text{meta}}_k \mid \mathcal{G}_k)],\; 1.0))$
    \ENDIF
\ENDFOR
\end{algorithmic}
\end{algorithm}


\subsubsection{Design Rationale}
\label{art1:sec:gapf_rationale}

Five design choices distinguish GAPF from standard ensemble methods: \textbf{(1)}~\emph{Topology awareness}: the GCN encodes the spatial graph, allowing the controller to learn that certain topologies favour one expert over another. \textbf{(2)}~\emph{Episode-level granularity}: selection occurs once per episode, avoiding instability from switching experts mid-episode where GRU-based agents have incompatible hidden state dynamics. \textbf{(3)}~\emph{Minimal overhead}: $1{,}473$ parameters vs.\ $38{,}983$ per expert; the GCN runs once per episode and the CAS decision is a single MLP forward pass. \textbf{(4)}~\emph{Frozen experts}: pre-trained policies are not fine-tuned, preserving validated performance and reducing the optimisation landscape to $1{,}473$ dimensions. \textbf{(5)}~\emph{Deployment feasibility}: each agent's GRU ($38{,}983$ parameters $\approx 38$\,KB in int8 quantisation) fits within typical WSN microcontroller RAM (nRF52840: 256\,KB; ESP32: 520\,KB), enabling on-device inference.


Here is the complete time-complexity analysis. 
\subsection{Time Complexity Analysis}
\label{art1:sec:complexity}

Let $n$ denote the number of sensor agents, $h = 64$ the RNN hidden dimension, $d = 16$ the GNN embedding dimension, and $T$ the number of steps per episode.

\paragraph{Per-step inference.}
At each routing step, only the frozen expert (QMIX or QTRAN) executes. Each agent $i$ runs: an input projection $\mathbf{x}_i = \text{ReLU}(\mathbf{W}_1 \mathbf{o}_i)$ in $\mathcal{O}(d_{\text{obs}} \cdot h)$, a GRU update with three $h \times h$ gates in $\mathcal{O}(h^2)$, and a Q-value output $q_i = \mathbf{W}_2 h_i^t$ in $\mathcal{O}(h \cdot n)$ (since the action space has dimension $n+1$, one per potential relay target). Across all $n$ agents:
\begin{equation}
\boxed{\mathcal{O}_{\text{step}} = \mathcal{O}(n^2 h + n h^2)}
\end{equation}
This is \textbf{identical for QMIX, QTRAN, and GAPF} at execution time---the mixer is not used during decentralised execution.

\paragraph{Episode-level GAPF overhead.}
GAPF adds a one-time computation before each episode: adjacency construction and normalisation in $\mathcal{O}(n^2)$, a two-layer GCN forward pass (sparse matrix--dense multiplications with weight matrices $\mathbf{W}_1 \in \mathbb{R}^{4 \times d}$, $\mathbf{W}_2 \in \mathbb{R}^{d \times d}$, and mean-pool readout) in $\mathcal{O}(n^2 d + n d^2)$, and the CAS MLP (two layers with hidden dimension $m = 64$) in $\mathcal{O}(dm)$. Since $d, m \ll n$:
\begin{equation}
\boxed{\mathcal{O}_{\text{GAPF-init}} = \mathcal{O}(n^2 d)}
\end{equation}

\paragraph{Total episode complexity.}
For an episode of $T$ steps: $\mathcal{O}_{\text{GAPF}} = \mathcal{O}(Tn^2 h + Tnh^2 + n^2 d)$. Since $Th \gg d$ (with $T = 100$, $h = 64$, $d = 16$), the GCN overhead is amortised and the dominant term $\mathcal{O}(Tn^2 h)$ is identical to standalone QMIX/QTRAN.

\begin{center}
\captionof{table}{Time complexity comparison. $n$: number of sensors, $T$: steps per episode, $h$: RNN hidden dimension (64), $d$: GNN embedding dimension (16), $N_p$: PSO particles (32), $I$: PSO iterations (50).}
\label{art1:tab:complexity}
\begin{tabular}{lcc}
\toprule
Algorithm & Per-Step & Per-Episode \\
\midrule
QMIX / QTRAN & $\mathcal{O}(n^2 h + n h^2)$ & $\mathcal{O}(T(n^2 h + n h^2))$ \\
QPSOFL  & $\mathcal{O}(N_p \cdot n^2)$ & $\mathcal{O}(T \cdot I \cdot N_p \cdot n^2)$ \\
\textbf{GAPF} & $\mathcal{O}(n^2 h + n h^2)$ & $\mathcal{O}(T(n^2 h + n h^2) + n^2 d)$ \\
\bottomrule
\end{tabular}
\end{center}

\noindent GAPF adds only $\mathcal{O}(n^2 d)$ per episode, where $d = 16 \ll h = 64$, making its per-step complexity identical to standalone QMIX/QTRAN. In contrast, QPSOFL's cost grows as $\mathcal{O}(T \cdot I \cdot N_p \cdot n^2)$ due to repeated swarm optimisation at every step, explaining its longer wall-clock times despite having no neural network.



\section{Results and Discussion}

Table~\ref{art1:table_setting} presents the custom Gym environment setting.

\begin{center}
\captionof{table}{Custom Gym Environment setting.}
\label{art1:table_setting}
\resizebox{0.6\textwidth}{!}{
\begin{tabular}{|>{\columncolor{gray!50}}c|>{\columncolor{gray!10}}c|}
\hline
\textbf{Key} & \textbf{Value} \\
\hline
$E_{\text{elec}}$ & $50 \times 10^{-9}$ J/bit \\
\hline
$E_{\text{amp}}$ & $100 \times 10^{-12}$ J/(bit$\cdot$m$^2$) \\
\hline
Packet size & 512 bit \\
\hline
Initial energy per sensor ($E_0$) & 1 J \\
\hline
Field bounds (LowerBound, UpperBound) & 0--100 m \\
\hline
Base station position & (50, 50) \\
\hline
Initial packets per sensor & 1 \\
\hline
Total timesteps / Max per episode & 200,000 / 100 \\
\hline
\end{tabular}
}
\end{center}

Table~\ref{art1:tab:hyperparameters} displays the training hyperparameters. QMIX and QTRAN are trained independently as separate expert policies; their parameters are then frozen, and only the lightweight GAPF meta-controller (${\sim}1{,}473$ parameters) is trained to dynamically select between them based on network topology.

\begin{center}
\captionof{table}{Training Hyperparameters for QMIX, QTRAN, and GAPF.}
\label{art1:tab:hyperparameters}
\resizebox{\textwidth}{!}{
\begin{tabular}{llccc}
\toprule
\textbf{Category} & \textbf{Parameter} & \textbf{QMIX} & \textbf{QTRAN} & \textbf{GAPF} \\
\midrule
\multirow{3}{*}{Optimization}
& Optimizer / LR       & Adam / 0.0003 & RMSprop / 0.0003 & Adam / 0.003 \\
& Grad clip / $\gamma$ & 10 / 0.99     & 10 / 0.99        & 1.0 / ---    \\
& Reward scale         & 0.01          & 0.01             & ---          \\
\midrule
\multirow{2}{*}{Exploration}
& Strategy             & $\varepsilon$-greedy & $\varepsilon$-greedy & REINFORCE \\
& $\varepsilon$: start$\to$end (anneal) & $0.5\to0.01$ (150K) & $0.5\to0.01$ (150K) & --- \\
\midrule
\multirow{2}{*}{Architecture}
& Agent / hidden dim   & GRU(64) / 64  & GRU(64) / 64  & GCN+MLP / 16+64 \\
& Mixing embed dim     & 64            & 64             & ---              \\
\midrule
\multirow{3}{*}{Training}
& Episodes / batch     & ${\sim}$2K / 32 & ${\sim}$2K / 32 & 2K / 1 (on-policy) \\
& Replay buffer        & 15,000       & 15,000        & ---              \\
& Target update        & Hard, every 30 ep. & Hard, every 30 ep. & EMA $\beta{=}0.95$ \\
\midrule
QTRAN-specific & $\lambda_{\text{opt}}$ / $\lambda_{\text{nopt}}$ & --- & 1.0 / 0.1 & --- \\
\midrule
Evaluation & Test episodes / eval $\varepsilon$ & 1,000 / 0.01 & 1,000 / 0.01 & 1,000 / --- \\
\bottomrule
\end{tabular}
}
\end{center}

Our evaluation includes QPSOFL~\cite{hu_energy_2024}, a metaheuristic baseline using Quantum PSO with Sobol-initialized particles and L'evy flight for cluster-head selection, followed by Mamdani-type Fuzzy Inference for relay selection.

\paragraph{Execution-time instrumentation.}
The framework records per-step wall-clock times (mean, p95, p99) as versioned \texttt{.npy} arrays. A standalone benchmark tool measures latency over 30 iterations (3 warmup), peak RAM via \texttt{tracemalloc}, analytical MACs/FLOPs, and estimated energy per inference ($\mu$J). Table~\ref{art1:tab:reviewer7_latency} reports results for the worst case ($N{=}70$).

\begin{center}
\captionof{table}{Inference metrics (Apple M-series CPU, 0.1,W reference). All algorithms achieve sub-millisecond latency.}
\label{art1:tab:reviewer7_latency}
\renewcommand{\arraystretch}{1.15}
\begin{tabular}{@{}lcccc@{}}
\toprule
\textbf{Algorithm} & \textbf{Latency p95 (ms)} & \textbf{Peak RAM (KB)} & \textbf{MACs} & \textbf{Energy ($\mu$J)} \\
\midrule
QMIX (agent + mixer)   & 0.27 & 2.98  & 2{,}539{,}904 & 27.2 \\
QTRAN (agent + mixer)   & 0.17 & 34.32 & 4{,}955{,}532 & 16.7 \\
GAPF (GCN, gateway)     & 0.03 & 4.12  & 181{,}408     & 3.0  \\
\midrule
Agent only (deployed)    & $<$1 & $\sim$2 & $\sim$29K  & $<$10 \\
\bottomrule
\end{tabular}
\end{center}

Even scaling by $50\times$ for a resource-constrained ARM Cortex-M4 at 80,MHz, the per-step decision time remains under 15,ms---well within real-time WSN budgets.

Seven scenarios (Table~\ref{art1:tab:scenarios}) systematically vary network density ($n$) and coverage radius ($R$), each evaluated over three random seeds (42, 123, 456).

\begin{center}
\captionof{table}{Overview of the seven evaluation scenarios.}
\label{art1:tab:scenarios}
\begin{tabular}{cccl}
\hline
\textbf{\#} & $n$ & $R$ (m) & \textbf{Rationale} \\
\hline
1 & 20  & 25 & Sparse, short range --- extreme connectivity challenge \\
2 & 30  & 25 & Small--medium, short range --- limited relay options \\
3 & 50  & 35 & Medium, standard range --- moderate density \\
4 & 70  & 35 & Baseline --- dense, standard range \\
5 & 70  & 45 & Dense, extended range --- more routing options \\
6 & 100 & 35 & Large-scale, standard range --- scalability test \\
7 & 100 & 50 & Large-scale, long range --- high connectivity stress \\
\hline
\end{tabular}
\end{center}

\subsubsection{Consolidated Results Across All Scenarios}

Table~\ref{art1:tab:consolidated_perf} presents key performance metrics for all seven scenarios. All values are mean $\pm$ std aggregated over three seeds on the 1{,}000-episode test set. PDR denotes packet delivery ratio, $E_{\text{tot}}$ total energy consumed (J), $\sigma_{E}$ standard deviation of remaining energy (lower is better balanced), and Throughput is packets delivered per step.

\begin{center}
\captionof{table}{Performance metrics across all scenarios (three seeds, test set). Best values per scenario in \textbf{bold}.}
\label{art1:tab:consolidated_perf}
\resizebox{\textwidth}{!}{
\begin{tabular}{cc|cccc|cccc}
\toprule
& & \multicolumn{4}{c|}{\textbf{PDR}} & \multicolumn{4}{c}{\textbf{$E_{\text{tot}}$ (J)}} \\
$n$ & $R$ & QMIX & QTRAN & QPSOFL & \textbf{GAPF} & QMIX & QTRAN & QPSOFL & \textbf{GAPF} \\
\midrule
20  & 25 & $.559\pm.257$ & $.558\pm.258$ & $.075\pm.107$ & $\mathbf{.753\pm.161}$ & $.093\pm.055$ & $.093\pm.055$ & $\mathbf{.052\pm.049}$ & $.075\pm.044$ \\
30  & 25 & $.568\pm.179$ & $.567\pm.174$ & $.053\pm.046$ & $\mathbf{.877\pm.159}$ & $.165\pm.054$ & $.166\pm.054$ & $\mathbf{.102\pm.050}$ & $.102\pm.070$ \\
50  & 35 & $.553\pm.248$ & $.552\pm.250$ & $.153\pm.092$ & $\mathbf{.990\pm.088}$ & $.423\pm.152$ & $.422\pm.153$ & $\mathbf{.113\pm.072}$ & $.202\pm.131$ \\
70  & 35 & $.545\pm.242$ & $.545\pm.243$ & $.122\pm.090$ & $\mathbf{.980\pm.119}$ & $.600\pm.204$ & $.603\pm.200$ & $\mathbf{.198\pm.116}$ & $.283\pm.213$ \\
70  & 45 & $.667\pm.246$ & $.659\pm.256$ & $.365\pm.092$ & $\mathbf{.980\pm.124}$ & $.687\pm.265$ & $.692\pm.274$ & $\mathbf{.040\pm.040}$ & $.334\pm.272$ \\
100 & 35 & $.464\pm.270$ & $.462\pm.261$ & $.099\pm.080$ & $\mathbf{.980\pm.124}$ & $1.010\pm.295$ & $1.017\pm.282$ & $\mathbf{.312\pm.171}$ & $.334\pm.272$ \\
100 & 50 & $.669\pm.266$ & $.684\pm.270$ & $.448\pm.057$ & $\mathbf{.980\pm.124}$ & $1.181\pm.413$ & $1.180\pm.412$ & $\mathbf{.026\pm.025}$ & $.334\pm.272$ \\
\bottomrule
\end{tabular}
}
\end{center}

\begin{center}
\captionof{table}{Energy balance ($\sigma_E$) and feasibility compliance across all scenarios.}
\label{art1:tab:consolidated_balance}
\resizebox{\textwidth}{!}{
\begin{tabular}{cc|cccc|cc|cc}
\toprule
& & \multicolumn{4}{c|}{$\sigma_E$ (lower = better)} & \multicolumn{2}{c|}{Memory $<1024$,MB?} & \multicolumn{2}{c}{$p_{99}<50$,ms?} \\
$n$ & $R$ & QMIX & QTRAN & QPSOFL & \textbf{GAPF} & QMIX/QTRAN & GAPF & QMIX/QTRAN & GAPF \\
\midrule
20  & 25 & .0075 & .0075 & .0056 & $\mathbf{.0043}$ & FAIL & PASS & PASS & PASS \\
30  & 25 & .0111 & .0112 & .0085 & $\mathbf{.0035}$ & FAIL & PASS & PASS & PASS \\
50  & 35 & .0264 & .0265 & .0081 & $\mathbf{.0039}$ & FAIL & PASS & PASS & PASS \\
70  & 35 & .0318 & .0319 & .0114 & $\mathbf{.0049}$ & FAIL & PASS & PASS & PASS \\
70  & 45 & .0379 & .0382 & $\mathbf{.0024}$ & .0061 & FAIL & PASS & PASS & PASS \\
100 & 35 & .0455 & .0457 & .0144 & $\mathbf{.0061}$ & FAIL & PASS & PASS & PASS \\
100 & 50 & .0559 & .0559 & $\mathbf{.0010}$ & .0061 & FAIL & PASS & PASS & PASS \\
\bottomrule
\end{tabular}
}
\end{center}

\subsubsection{Key Findings}

Five principal findings emerge from the seven scenarios:

\paragraph{(F1) GAPF consistently achieves the highest PDR, with the advantage growing in sparser networks.}
GAPF's relative PDR improvement over the best standalone MARL baseline scales from $+35\%$ ($n{=}20$) to $+54\%$ ($n{=}30$), $+79\%$ ($n{=}50$), and peaks at $+111\%$ ($n{=}100$, $R{=}35$). Increasing the coverage radius narrows the gap (e.g., $+80\%$ at $(70,35)$ vs $+47\%$ at $(70,45)$; $+111\%$ at $(100,35)$ vs $+43\%$ at $(100,50)$), confirming that GAPF's topology-aware policy fusion is most valuable where routing is hardest---sparse networks with limited relay options.

\paragraph{(F2) GAPF achieves $2$--$3.5\times$ lower energy consumption and $3$--$9\times$ better energy balance than QMIX/QTRAN.}
While QPSOFL uses the least total energy in most scenarios, its near-zero delivery renders this meaningless. Among algorithms with viable PDR ($>50\%$), GAPF consistently uses the least energy. The energy balance ratio ($\sigma_E^{\text{QMIX}}/\sigma_E^{\text{GAPF}}$) widens from $1.7\times$ ($n{=}20$) to $9.2\times$ ($n{=}100$, $R{=}50$), confirming that GAPF's GCN-based load distribution scales with network size.

\paragraph{(F3) QMIX and QTRAN converge to statistically indistinguishable performance.}
Across all scenarios, QMIX and QTRAN yield nearly identical PDR (within 1--2\%), energy, and latency. This validates GAPF's design: the performance gain stems from the GCN-based topology-aware controller and policy fusion, not from expert selection alone.

\paragraph{(F4) QPSOFL collapses at scale; low energy reflects non-delivery.}
QPSOFL's PDR degrades from $7.5\%$ ($n{=}20$) to below $10\%$ at $n{=}100$, $R{=}35$. Wider coverage radii partially recover delivery ($44.8\%$ at $n{=}100$, $R{=}50$), but performance remains unacceptable. Its excellent energy balance at high $R$ (e.g., $\sigma_E{=}0.001$ at $(100,50)$) is an artifact of minimal packet transmission rather than intelligent routing.

\paragraph{(F5) GAPF is the only algorithm passing all feasibility constraints.}
QMIX/QTRAN \textbf{fail} the 1{,}024,MB memory constraint in every scenario (peak memory $3.2$--$13.8$,GB), while GAPF stays within budget ($684$--$951$,MB). All algorithms pass the $p_{99}{<}50$,ms latency constraint. GAPF's memory footprint at $n{=}100$ ($951$,MB) leaves only $7\%$ headroom, warranting attention for even larger deployments.

\subsection{Deployment Feasibility on Resource-Constrained Nodes}
\label{art1:sec:deployment_feasibility}

The CTDE paradigm~\cite{kraemer2016multi,lowe2017multi} decouples training from deployment complexity. At inference time, the heavy training-only components---value-decomposition mixers, reward scalarization network, ES loop---are \emph{discarded entirely}.

\begin{center}
\captionof{table}{Component deployment map under CTDE. Only the per-sensor agent and GAPF controller are active at inference.}
\label{art1:tab:deployment_components}
\renewcommand{\arraystretch}{1.15}
\resizebox{\textwidth}{!}{
\begin{tabular}{@{}lccc@{}}
\toprule
\textbf{Component} & \textbf{Runs on} & \textbf{Deployed?} & \textbf{Size} \\
\midrule
QMIX/QTRAN mixer, Q$\cdot$K Attention, ES loop & Training server & No & --- \\
\midrule
GAPF meta-controller (GCN + CAS) & Sink/gateway & Yes (1$\times$/episode) & $\sim$1{,}473 params ($\sim$6,KB) \\
Per-sensor agent (GRU) & Each sensor & Yes (1$\times$/step) & $\sim$25K params ($\sim$100,KB) \\
\bottomrule
\end{tabular}
}
\end{center}

Each sensor executes a single-layer GRU agent:
\begin{equation}
a_i = \argmax_{j}; \mathbf{W}_2 ;\mathrm{GRU}!\bigl(\mathrm{ReLU}(\mathbf{W}_1 \mathbf{o}_i + \mathbf{b}1),; \mathbf{h}{i}^{(t-1)}\bigr) + \mathbf{b}_2
\label{art1:eq:agent_forward}
\end{equation}
where $\mathbf{o}_i \in \mathbb{R}^5$ is the local observation (remaining energy, consumed energy, $x$-$y$ position, packet count), $\mathbf{h}_i^{(t-1)} \in \mathbb{R}^{64}$ is the recurrent hidden state, and $|\mathcal{A}|{=}N{+}1$. A single forward pass requires ${\sim}29$K MACs (${\sim}320$ for the input layer, ${\sim}24{,}576$ for the GRU cell, ${\sim}4{,}544$ for the output layer at $n{=}70$), which completes in $<1$,ms even on an ARM Cortex-M4 at 80,MHz. The full model fits in ${\sim}100$,KB of flash, well within the capacity of commercial WSN platforms such as the STM32L4 (1,MB flash, 320,KB SRAM) or nRF52840 (1,MB flash, 256,KB RAM).

\subsection{Limitations}
\label{art1:sec:limitations}

Despite GAPF's strong empirical performance, several limitations warrant discussion:

\begin{enumerate}
    \item \textbf{Simulation-only validation.} All experiments are conducted in a custom OpenAI Gym simulator. While we provide deployment feasibility analysis (Section~\ref{art1:sec:deployment_feasibility}) showing that the deployed agent requires only ${\sim}100$ lines of C and $<30$\,\textmu J per inference on a Cortex-M4 MCU, no hardware-in-the-loop experiments have been performed. Real-world factors such as channel fading, packet collisions, and clock drift are not modelled.
    
    \item \textbf{Static topology assumption.} GAPF computes the graph embedding once per episode and selects a fixed expert for the entire episode. In scenarios with mobile sensors or rapidly changing channel conditions, step-level re-evaluation (using the Fusion controller) may be necessary, at the cost of increased computational overhead.
    
    \item \textbf{Two-expert limitation.} The current CAS controller produces a binary selection (QMIX vs QTRAN). Extending to $k > 2$ experts (e.g., incorporating DQN, MAPPO, or QPLEX) would require a softmax-based selection mechanism and raises questions about diminishing returns as expert diversity increases.
    
    \item \textbf{Fixed episode length.} All scenarios use $T = 100$ steps per episode. Real WSN deployments operate continuously with time-varying traffic patterns and gradual energy depletion over days or weeks. The episodic formulation does not capture long-horizon energy management.
    
    \item \textbf{Memory constraint tightness.} At $n{=}100$, GAPF's memory footprint reaches 951\,MB---within the 1\,GB constraint but with only 7\% headroom. Scaling beyond $n{=}100$ may require model compression techniques (pruning, quantisation) to remain feasible on resource-constrained gateways.
    
    \item \textbf{Homogeneous sensor assumption.} All sensors are identical in hardware capability and initial energy. Heterogeneous deployments with mixed sensor types, varying transmission power, or asymmetric energy budgets are not evaluated.
\end{enumerate}

\FloatBarrier



\section{Conclusion}
\label{art1:sec:conclusion}

This paper introduced \textbf{Graph-Attentive Policy Fusion (GAPF)}, a lightweight meta-learning framework for energy-efficient multi-hop routing in Wireless Sensor Networks. GAPF fuses frozen MARL expert policies (QMIX, QTRAN) through a two-layer Graph Convolutional Network that encodes the sensor topology into a compact 16-dimensional embedding, enabling a Contextual Algorithm Selection controller to choose the best expert per deployment with only 1{,}473 trainable parameters.

\subsection{Contributions and Impact}

Extensive evaluation across 7 scenarios (20--100 sensors, 25--50\,m coverage radius, 3 seeds, 21{,}000 test episodes) establishes three principal contributions:

\begin{enumerate}
    \item \textbf{Near-perfect, scale-invariant delivery.} GAPF achieves $\geq$98\% packet delivery ratio from 20 to 100 sensors, with +43\% to +111\% relative improvement over standalone MARL. This reliability is critical for \emph{life-saving responsiveness} in disaster monitoring: in earthquake, flood, or wildfire scenarios, every undelivered sensor reading represents a potential blind spot in situational awareness that could cost lives.
    
    \item \textbf{Energy-coherent routing for green AI.} GAPF consumes 2--3.5$\times$ less energy than QMIX/QTRAN while delivering more data, and maintains 3--9$\times$ better energy balance across sensors. In battery-powered WSN deployments, this translates directly to \emph{extended network lifetime}---reducing the frequency of battery replacement, which in disaster zones may be infeasible. The environmental impact is twofold: fewer batteries consumed (reducing e-waste) and longer monitoring coverage per deployment (reducing the carbon footprint of redeployment logistics).
    
    \item \textbf{Feasible deployment on constrained hardware.} GAPF passes all feasibility constraints ($p_{99} < 50$\,ms, peak memory $< 1{,}024$\,MB) across all scenarios, while QMIX/QTRAN exceed the memory budget by up to $13\times$. The Centralised Training with Decentralised Execution (CTDE) paradigm ensures that only a lightweight GRU agent (${\sim}3$\,KB weights, ${\sim}27$\,\textmu J per inference) runs on each sensor node, making GAPF compatible with commercial microcontrollers (ARM Cortex-M4, 256\,KB SRAM) without requiring cloud connectivity during operation.
\end{enumerate}

\subsection{Broader Impact}

Beyond WSN routing, GAPF demonstrates a general principle: \emph{lightweight policy fusion can outperform heavyweight monolithic models} in multi-agent systems. The approach of encoding structural context via GNNs and selecting among frozen specialist policies is applicable to other domains where topology matters---vehicular ad-hoc networks, drone swarm coordination, and smart grid load balancing. By keeping the meta-controller small (1{,}473 parameters vs millions in end-to-end approaches), GAPF aligns with the growing imperative for \emph{responsible AI}: models that are interpretable (binary expert selection is auditable), efficient (minimal carbon footprint for training and inference), and deployable without expensive GPU infrastructure.

\subsection{Future Work}

We identify five directions for future research, ordered by expected impact:

\begin{enumerate}
    \item \textbf{Hardware-in-the-loop validation.} Deploying GAPF on physical WSN testbeds (e.g., TelosB or nRF52840 motes) with realistic channel models, packet collisions, and environmental interference. Weight extraction to embedded C (Section~\ref{art1:sec:deployment_feasibility}) provides a concrete deployment path; the remaining gap is empirical validation under real RF conditions.
    
    \item \textbf{Cross-domain disaster monitoring.} Adapting the framework to heterogeneous disaster scenarios---flood monitoring (mobile water-level sensors), wildfire detection (temperature/smoke gradients with rapidly evolving topology), and post-earthquake structural health monitoring (static sensors with intermittent connectivity). Each domain introduces unique topology dynamics that GAPF's GCN encoder is architecturally suited to capture.
    
    \item \textbf{Model compression for extreme scale.} At $n{=}100$, GAPF's gateway memory reaches 951\,MB. Techniques such as 8-bit quantisation, structured pruning of expert policies, and knowledge distillation into a single lightweight policy could extend feasibility to $n > 200$ while reducing inference energy further---a direct contribution to \emph{green AI} in edge computing.
    
    \item \textbf{Dynamic expert fusion.} Extending the CAS controller to step-level fusion (blending Q-values from multiple experts with learned mixing weights) rather than episode-level binary selection. This would enable GAPF to adapt to \emph{within-episode} topology changes caused by sensor failures or energy depletion, at the cost of additional per-step computation.
    
    \item \textbf{Multi-expert generalisation.} Incorporating additional MARL algorithms (QPLEX, MAPPO, DQN) as experts and investigating whether expert diversity yields diminishing or compounding returns. A softmax-based selection over $k$ experts with the same GCN encoder would test the limits of the policy fusion paradigm.
\end{enumerate}

\medskip
\noindent In summary, GAPF bridges the gap between the theoretical promise of multi-agent reinforcement learning and the practical constraints of real-world sensor networks. By combining graph-structural intelligence with policy fusion, it achieves what neither MARL nor metaheuristic approaches accomplish alone: reliable, energy-efficient, and deployable routing at scale---a step toward AI systems that are not only performant but also responsible, sustainable, and directly impactful for public safety.


\vspace{0.5cm}
\noindent Les hyperparamètres complets, la preuve formelle de monotonicité, les résultats détaillés par scénario, l'analyse de faisabilité de déploiement et un tableau comparatif des approches de routage récentes (2024--2025) sont fournis à l'Annexe~\ref{ann:gapf-suppl}.
